% Part: first-order-logic
% Chapter: syntax-and-semantics
% Section: first-order-languages

\documentclass[../../../include/open-logic-section]{subfiles}

\begin{document}

\olfileid{fol}{syn}{fol}

\olsection{زبان‌های مرتبهٔ اول}


عبارت‌های منطق مرتبهٔ اول از واژگانی پایه ساخته می‌شوند که
\emph{!!{variable}s}، \emph{!!{constant}s}،
\emph{!!{predicate}s} و گاه \emph{!!{function}s} را در بر می‌گیرد.
از این‌ها، همراه با رابط‌های منطقی، سورها و نشانه‌های سجاوندی‌ای چون
پرانتز و ویرگول، \emph{ترم‌ها} و \emph{!!{formula}s} ساخته می‌شوند.

\begin{explain}
به‌طور غیرصوری، !!{predicate}s نام ویژگی‌ها و رابطه‌ها هستند،
!!{constant}s نام اشیای منفردند و !!{function}s نام نگاشت‌ها. این‌ها، به‌جز
!!{identity}~$\eq$، \emph{نمادهای غیرمنطقی}اند و در کنار هم زبانی را
می‌سازند. هر زبان مرتبهٔ اول~$\Lang L$ با نمادهای غیرمنطقی‌اش تعیین می‌شود.
در کلی‌ترین حالت، $\Lang L$ از هر نوع بی‌نهایت نماد در بر دارد.
\end{explain}

در حالت کلی، در منطق مرتبهٔ اول از نمادهای زیر استفاده می‌کنیم:

\begin{enumerate}
\item نمادهای منطقی
\begin{enumerate}
\item رابط‌های منطقی:
  \startycommalist
  \iftag{prvNot}{\ycomma $\lnot$ (نفی)}{}%
  \iftag{prvAnd}{\ycomma $\land$ (عطف)}{}%
  \iftag{prvOr}{\ycomma $\lor$ (فصل)}{}%
  \iftag{prvIf}{\ycomma $\lif$ (!!{conditional})}{}%
  \iftag{prvIff}{\ycomma $\liff$ (!!{biconditional})}{}%
  \iftag{prvAll}{\ycomma $\lforall$ (سور کلی)}{}%
  \iftag{prvEx}{\ycomma $\lexists$ (سور وجودی)}{}.
\tagitem{prvFalse}{ثابت گزاره‌ایِ !!{falsity}~$\lfalse$.}{}
\tagitem{prvTrue}{ثابت گزاره‌ایِ !!{truth}~$\ltrue$.}{}
\item !!{identity} دوجایگاهی~$\eq$.
\item مجموعه‌ای !!{denumerable} از !!{variable}s: $\Obj v_0$، $\Obj v_1$، $\Obj
  v_2$، \dots
\end{enumerate}
\item نمادهای غیرمنطقی، که \emph{زبان استاندارد} منطق مرتبهٔ اول را
تشکیل می‌دهند
\begin{enumerate}
\item برای هر $n>0$، مجموعه‌ای !!{denumerable} از !!{predicate}sی
  $n$جایگاهی: $\Obj A^n_0$، $\Obj A^n_1$، $\Obj A^n_2$، \dots
\item مجموعه‌ای !!{denumerable} از !!{constant}s: $\Obj c_0$، $\Obj c_1$، $\Obj
  c_2$، \dots.
\item برای هر $n>0$، مجموعه‌ای !!{denumerable} از !!{function}sی
  $n$جایگاهی: $\Obj f^n_0$، $\Obj f^n_1$، $\Obj f^n_2$، \dots
\end{enumerate}
\item نشانه‌های سجاوندی: ( و )، و نیز ویرگول.
\end{enumerate}

بیشتر تعریف‌ها و نتایج خود را برای زبان استاندارد کامل منطق مرتبهٔ اول
صورت‌بندی خواهیم کرد. بااین‌حال، بسته به کاربرد، ممکن است زبان را تنها به
چند !!{predicate}، !!{constant} و !!{function} محدود کنیم.

\begin{ex}
زبان حساب~$\Lang L_A$ یک !!{predicate} دوجایگاهی~$<$، یک
!!{constant}~$\Obj 0$، یک !!{function} یک‌جایگاهی~$\prime$ و دو
!!{function} دوجایگاهی~$+$ و~$\times$ دارد.
\end{ex}

\begin{ex}
زبان نظریهٔ مجموعه‌ها~$\Lang L_Z$ تنها !!{predicate} دوجایگاهی~$\in$ را
دارد.
\end{ex}

\begin{ex}
زبان ترتیب‌ها~$\Lang L_\le$ تنها !!{predicate} دوجایگاهی~$\le$ را دارد.
\end{ex}

باز هم این‌ها قراردادند: به‌طور رسمی، صرفاً نام‌های مستعارند؛ برای نمونه،
$<$، $\in$ و $\le$ نام‌های مستعار $\Obj A^2_0$، $\Obj 0$ نام مستعار
$\Obj c_0$، $\prime$ نام مستعار $\Obj f^1_0$، $+$ نام مستعار $\Obj f^2_0$
و $\times$ نام مستعار $\Obj f^2_1$ است.

\iftag{defNot,defOr,defAnd,defIf,defIff,defTrue,defFalse,defEx,defAll}{%
افزون بر رابط‌ها و
\iftag{notprvEx,notprvAll}{سور}{سورهای} اولیه‌ای که در بالا معرفی شدند،
از نمادهای \emph{تعریف‌شدهٔ} زیر نیز استفاده می‌کنیم:
\startycommalist
  \iftag{defNot}{\ycomma $\lnot$ (نفی)}{}%
  \iftag{defAnd}{\ycomma $\land$ (عطف)}{}%
  \iftag{defOr}{\ycomma $\lor$ (فصل)}{}%
  \iftag{defIf}{\ycomma $\lif$ (!!{conditional})}{}%
  \iftag{defIff}{\ycomma $\liff$ (!!{biconditional})}{}%
  \iftag{defAll}{\ycomma $\lforall$ (سور کلی)}{}%
  \iftag{defEx}{\ycomma $\lexists$ (سور وجودی)}{}%
  \iftag{defFalse}{\ycomma !!{falsity}~$\lfalse$}{}%
  \iftag{defTrue}{\ycomma !!{truth}~$\ltrue$}{}.
}{}

\begin{tagblock}{defNot,defOr,defAnd,defIf,defIff,defTrue,defFalse,defEx,defAll}
\begin{explain}
نماد تعریف‌شده رسماً جزئی از زبان نیست، بلکه به‌عنوان اختصاری غیرصوری
معرفی می‌شود: چنین نمادی به ما امکان می‌دهد فرمول‌هایی را کوتاه کنیم که
اگر فقط از نمادهای اولیه استفاده می‌کردیم، بسیار طولانی می‌شدند. این مزیت
آشکار است؛ اما مزیت بزرگ‌تر آن است که برهان‌ها کوتاه‌تر می‌شوند. اگر نمادی
اولیه باشد، باید در برهان‌ها جداگانه بررسی شود. بنابراین هرچه شمار نمادهای
اولیه بیشتر باشد، برهان‌های ما طولانی‌تر خواهند بود.
\end{explain}
\end{tagblock}

% Alternate symbols

\begin{intro}
ممکن است با اصطلاح‌ها و نمادهایی متفاوت از آنچه در بالا به کار بردیم آشنا
باشید. متون منطق (و مدرسان) معمولاً $\sim$، $\neg$ یا~!\ را برای «نفی» و
$\wedge$، $\cdot$ یا $\&$ را برای «عطف» به کار می‌برند. نمادهای رایج برای
«شرطی» یا «استلزام» عبارت‌اند از $\rightarrow$، $\Rightarrow$ و $\supset$.
\iftag{prvIff,defIff}{نمادهای «دوشرطی»، «استلزام دوسویه» یا «هم‌ارزی
  (مادی)» عبارت‌اند از $\leftrightarrow$،
  $\Leftrightarrow$ و $\equiv$.}{}
\iftag{prvFalse,defFalse}{نماد $\lfalse$ را به نام‌های «کذب»، «فالسوم»،
  «محال» یا «پایین» نیز می‌خوانند.}{}
\iftag{prvTrue,defTrue}{نماد $\ltrue$ را به نام‌های «صدق»، «وروم» یا
  «بالا» نیز می‌خوانند.}{}

بنا بر قرارداد، برای !!{constant}s (که گاه نام نیز خوانده می‌شوند) از حروف
کوچک آغاز الفبای لاتین (مانند $a$، $b$ و $c$) و برای !!{variable}s از حروف
کوچک پایان آن (مانند $x$، $y$ و $z$) استفاده می‌شود. سورها با
!!{variable}s، مانند $x$، ترکیب می‌شوند؛ صورت‌های نگارشی متفاوت برای سور
کلی عبارت‌اند از $\forall x$، $(\forall x)$، $(x)$، $\Pi x$ و
$\bigwedge_x$، و برای سور وجودی عبارت‌اند از $\exists x$، $(\exists
x)$، $(Ex)$، $\Sigma x$ و $\bigvee_x$.
\end{intro}

\begin{explain}
می‌توانیم همهٔ عملگرهای گزاره‌ای و هر دو سور را نمادهای اولیهٔ زبان به شمار
آوریم. یا می‌توانیم مجموعهٔ کوچک‌تری از نمادهای اولیه برگزینیم و دیگر
!!{operator}s را تعریف‌شده بگیریم. مجموعه‌های «تام از نظر تابع ارزش» از
عملگرهای بولی شامل $\{ \lnot, \lor \}$، $\{ \lnot, \land \}$ و
$\{ \lnot, \lif\}$ هستند---هر یک از این‌ها را می‌توان با یکی از دو سور
ترکیب کرد تا زبانی مرتبهٔ اول و تام از نظر قدرت بیان به دست آید.

ممکن است با دو !!{operator} دیگر نیز آشنا باشید: خط شفر~$|$ (به نام هنری
شفر) و پیکان پیرس~$\downarrow$، که خنجر کواین نیز نامیده می‌شود. این دو
عملگر، با خوانش معمول «نَند» و «نُر» (به‌ترتیب)، هر یک به‌تنهایی از نظر
تابع ارزش تام‌اند.
\end{explain}

\end{document}
